% noinspection SpellCheckingInspection


\documentclass[9pt,a4paper,twoside]{article}
\usepackage[utf8]{inputenc}
\usepackage[T1]{fontenc}
\usepackage[polish,latin]{babel}
\usepackage{geometry}
\usepackage{multicol}
\usepackage{lettrine}
\usepackage{xcolor}
\usepackage{titlesec}

% Ustawienia marginesów dla mszalika
\geometry{
    a4paper,
    top=15mm,
    bottom=15mm,
    left=15mm,
    right=15mm,
    footskip=10mm
}

% Kolory liturgiczne
\definecolor{rubricred}{RGB}{170, 0, 0}
\newcommand{\rubric}[1]{\textcolor{rubricred}{\textit{#1}}}

% Formatowanie nagłówków
\titleformat{\section}{\centering\normalfont\large\bfseries\color{rubricred}}{}{0pt}{}
\titleformat{\subsection}{\centering\normalfont\normalsize\bfseries\color{rubricred}}{}{0pt}{}

\begin{document}

% --- STRONA 1 (Nagłówek) ---
\begin{center}
    \vspace*{1cm}
    {\huge \textbf{PIERWSZA}} \\
    {\huge \textbf{KOMUNIA}} \\
    {\huge \textbf{ŚWIĘTA}} \\
    \vspace{5mm}
    {\large Warszawa-Radość/Józefów} \\
    \vspace{2mm}
    {\large 16 maja 2026 r.}
\end{center}

\newpage

% --- STRONA 3 (Początek Mszy) ---
\setcounter{page}{3}
\section*{ORDO MISSÆ}
\begin{center}
    \small \textit{16 maja, św. Andrzeja Boboli, szaty białe, I kl.}
\end{center}

\subsection*{MSZA KATECHUMENÓW}

\begin{center}
    \rubric{Pokropienie wodą święconą}
\end{center}
\noindent \footnotesize \rubric{W niedzielę kapłan kropi wiernych wodą święconą. Obrzęd ten przypomina nam, że jesteśmy odrodzeni wodą chrztu świętego. Śpiewamy Asperges me (s. 50)...}

\vspace{3mm}
\subsection*{Modlitwy u stopni ołtarza}
\noindent \footnotesize \rubric{Kapłan, stanąwszy u stopni ołtarza, żegna się znakiem krzyża i odmawia antyfonę:}

\begin{multicols}{2}
\selectlanguage{latin}
\noindent \lettrine[lines=2,findent=2pt,nindent=0pt]{\textcolor{rubricred}{I}}{N} nómine Patris, et Fílii, et Spíritus Sancti. Amen. \\
\textbf{S.} Introíbo ad altáre Dei. \\
\textbf{M.} Ad Deum, qui lætíficat iuventútem meam.

\columnbreak
\selectlanguage{polish}
\noindent \textbf{K.} W imię Ojca i Syna, i Ducha Świętego. Amen. \\
\textbf{K.} Przystąpię do ołtarza Bożego. \\
\textbf{M.} Do Boga, który jest weselem moim od młodości. [cite: 11, 12]
\end{multicols}

\vspace{-2mm}
\subsection*{Psalm 42}
\noindent \footnotesize \rubric{Opuszcza się we mszach żałobnych i w okresie Męki Pańskiej...} [cite: 13]

\begin{multicols}{2}
\selectlanguage{latin}
\noindent \textbf{S.} Iúdica me, Deus, et discérne causam meam de gente non sancta: ab hómine iníquo et dolóso érue me. \\
\textbf{M.} Quia tu es, Deus, fortitúdo mea: quare me repulísti, et quare tristis incédo, dum afflígit me inimícus? [cite: 14, 15]

\columnbreak
\selectlanguage{polish}
\noindent \textbf{K.} Bądź mi sędzią, o Boże, i rozsądź sprawę moją z narodem bezbożnym; \\
\textbf{K.} wybaw mnie od człowieka niedobrego i fałszywego. \\
\textbf{M.} Wszak Ty jesteś, o Boże, mocą moją; czemu mnie odrzucasz i czemu smutny chodzę, gdy nieprzyjaciel mnie nęka? [cite: 19, 20, 21]
\end{multicols}

\vspace{5mm}
\subsection*{Kyrie - błagalne wołanie}
\noindent \footnotesize \rubric{Kyrie jest to krótka litania pochodząca z liturgii greckiej...} [cite: 66, 67]

\begin{center}
\small
\begin{tabular}{llcll}
\selectlanguage{latin} \textbf{S.} KÝRIE, eléison. & & \textbf{\textcolor{rubricred}{OJCIEC}} & & \selectlanguage{polish} \textbf{K.} Panie, zmiłuj się. \\
\selectlanguage{latin} \textbf{M.} Kýrie, eléison. & & & & \selectlanguage{polish} \textbf{M.} Panie, zmiłuj się. \\
\hline
\selectlanguage{latin} \textbf{S.} Christe, eléison. & & \textbf{\textcolor{rubricred}{SYN}} & & \selectlanguage{polish} \textbf{K.} Chryste, zmiłuj się. \\
\selectlanguage{latin} \textbf{M.} Christe, eléison. & & & & \selectlanguage{polish} \textbf{M.} Chryste, zmiłuj się. \\
\end{tabular}
\end{center}

\end{document}